When a compact Calabi--Yau threefold develops singularities, M-theory on it acquires light
states, and the vacuum moduli space of the five-dimensional supergravity acquires a new branch
along which these states condense. With gravity decoupled, the local structure of such a Higgs
branch is a hyperk\"ahler quotient computed from the charges of the light states. In a compactification on a compact threefold the Planck mass is finite, the hypermultiplet scalars take values in a
quaternionic K\"ahler manifold rather than a hyperk\"ahler one~\cite{Galicki:1988}, and the
metric of that manifold is not known~\cite{Mohaupt:2004de}. This paper asks what the local space
of supersymmetric vacua is at finite Planck mass, and shows that for a class of transitions it
can be determined without knowing the metric.

We use the following terms. A \emph{root} of the Higgs branch is a point at which it meets the
Coulomb branch, that is, where the charged hypermultiplets become massless. If near a root the
Higgs branch is locally a product (vector moduli)$\,\times\,$(neutral
hypermultiplets)$\,\times\,S$ with $S$ a cone, we call the germ of $S$ at its vertex the
\emph{transverse type} of the root. The \emph{rigid model} of a root is the flat hyperk\"ahler
quotient of the light hypermultiplets by the gauge group that acts on them, with gravity and all
other fields ignored. In the language of stratified spaces the transverse type is the transverse
slice to the stratum through the root. Finally, a \emph{Wilsonian} effective action is one for
the modes below a fixed mass scale that keeps the light charged states as fields, as opposed to
the moduli-space action obtained by integrating them out.

\paragraph{The example.}
The threefold $X_9$ of~\cite{W26global,W26quantum}, recalled in Section~\ref{subsec:x9}, has
two ordinary double points and one cone over the degree-seven del Pezzo surface
$S_7=\operatorname{Bl}_2\mathbb P^2$; its crepant resolution $\wideX_9$ has Hodge numbers
$(h^{1,1},h^{2,1})=(6,122)$ and its smoothing $X_{9,t}$ has $(3,123)$. Physically, the two nodes
contribute two charged hypermultiplets and the del Pezzo cone the rank-one $E_2$ theory. This is
the strongly coupled five-dimensional superconformal theory, without a Lagrangian description,
obtained from M-theory on a local del Pezzo surface of degree seven; in the same series $E_1$
corresponds to $\mathbb P^1\times\mathbb P^1$ and $E_0$ to $\mathbb P^2$. Write $C_1,C_2$ for the exceptional
curves of small resolutions of the two nodes, $u_1,u_2$ for their smoothing parameters
($xy-zw=u_i$) and $s$ for the deformation parameter of the del Pezzo cone
(Section~\ref{subsec:x9}). The three sectors are charged under the same two bulk vector multiplets,
and every global first-order deformation of $X_9$ satisfies $u_1+s=u_2+s=0$, so none of them can be Higgsed
alone: every first-order smoothing of $X_9$ smooths all three singular points at once. The rigid model of
this transition has transverse type
\begin{equation}
 \CC^2/\ZZ_4,\qquad \CC[X,Y,s]/(XY-s^4),
 \label{eq:intro-A3}
\end{equation}
an $A_3$ surface singularity~\cite[eq.~(3.10)]{W26quantum}. The coupling of the three sectors
is a gravitational effect. The inverse gauge couplings of the vector multiplets are the Hodge
norms of the corresponding two-forms. Along any path on which the volume of the threefold grows
while the areas of $C_1$, $C_2$ and of all curves in the del Pezzo surface stay bounded, some
combination of the classes that couple the sectors acquires a divergent norm and its canonically
normalised coupling to the local fields vanishes. The joint constraint is then never retained; at
most $u_1=u_2$ survives~\cite[Sect.~4, Thm.~4.3]{W26quantum}. The question is therefore only
meaningful at finite Planck mass.

\paragraph{Results.}
First, we prove a local normal form for quaternionic K\"ahler quotients
(Theorem~\ref{thm:normal-form}). Let a compact connected group $G$ act isometrically on a
quaternionic K\"ahler manifold $M$ of dimension at least eight and nonzero scalar curvature, which
need not be complete, and
let $p$ lie on a component $N$ of the fixed set on which the moment map vanishes. If $G$ acts
locally freely on the nonzero part of the zero set of the flat hyperk\"ahler moment map $\mu_0$ of
the normal isotropy representation $V$, then near $[p]$ the quotient is homeomorphic, compatibly
with its stratification, to $N\times V/\!/\!/G$, where $V/\!/\!/G=\mu_0^{-1}(0)/G$ is the flat
hyperk\"ahler quotient. If the action there is free, the metric
tangent cone is the flat product, and the metric agrees with it up to relative corrections of
order $r^2$ at distance $r$ from $[p]$ (Proposition~\ref{prop:tangent-cone}); the rate is
attained. The statement holds for either sign of the scalar curvature, so it applies to positive
quaternionic K\"ahler manifolds as well. Behind this lies an
elementary observation (Lemma~\ref{lem:isotropy} and Remark~\ref{rem:no-fi}): the moment map of
a quaternionic K\"ahler manifold has no constants of integration, and at a fixed point it is
determined by the isotropy representation. A hyperk\"ahler quotient can be resolved or deformed
by moving the level of the moment map, a Fayet--Iliopoulos parameter; at finite Planck mass no
such parameter exists, and a nonzero value of the moment map at a fixed point is the same thing
as a gauging of (a subgroup of) the $SU(2)$ R-symmetry there. The normal form and the
tangent-cone estimate apply, at the level $\mu(p)$, to hyperk\"ahler and K\"ahler quotients at a fixed point
(Remark~\ref{rem:hk}). 

When the flat quotient $V/\!/\!/G$ is four-dimensional and the action is free off the origin, we also identify the correction of order
$r^2$ transverse to $N$. Suppose that $V$ contains a quaternionic line that lies in $\mu_0^{-1}(0)$ and
is orthogonal to the $G$-orbits through its points, as it does for the torus actions on $\HH^k$ with
one relation among the weights considered below. That line maps onto the flat quotient and gives
an orbifold chart of the transverse slice. The metric in this chart is not known to be twice
differentiable at the vertex, but it has an expansion to second order, and we call the curvature at the
vertex of its second-order truncation the \emph{leading curvature}. It is the curvature of $M$ at $p$ restricted to that
line (Proposition~\ref{prop:universal}): the part of the curvature fixed by supergravity
contributes constant curvature $4\nu$, where $\nu$ is the reduced scalar curvature of $M$ (with $\nu=1$ for
$\HH\mathbb P^n$ of sectional curvature in $[1,4]$), and the hyperk\"ahler part of the curvature contributes an
anti-self-dual Weyl tensor, with respect to the orientation in which the K\"ahler forms $\omega_a$ are
self-dual (self-dual in the opposite orientation customary in four-dimensional quaternionic K\"ahler
geometry).

Second, the physical statement holds for every face with one relation
(Theorem~\ref{thm:general}). Suppose that on an open face $F$ of the nef cone of a compact
Calabi--Yau threefold $\widehat X$ exactly $k$ disjoint $(-1,-1)$-curves have zero area, and that
their classes satisfy a single primitive relation $\sum_ia_i[C_i]=0$ with every $a_i\neq0$, that $F$ is
open in the annihilator of their classes, and that the gauge kinetic matrix of the vector multiplets is
positive definite on $F$ (Section~\ref{subsec:setting}).
Assume the regularity hypothesis of Section~\ref{subsec:hypotheses}, whose main content is that
there is a smooth Wilsonian hypermultiplet manifold containing the $k$ light states and that the
Coulomb branch through the root consists of supersymmetric vacua. Then the Higgs branch at the
root has transverse type $\CC^2/\ZZ_m$ with $m=\sum_i|a_i|$, the same as its rigid model, and
the metric of its hypermultiplet factor is the flat orbifold metric of $\CC^2/\ZZ_m$ times the flat
metric on the tangent space of $N$, up to relative corrections
of order $r^2$, where $r$ is the distance from the root. For $k$ nodes whose vanishing three-spheres are homologous up to sign, $m=k$.

Third, $X_9$ has such a face (Section~\ref{sec:face}). On a three-dimensional open face $F$ of
the nef cone of $\wideX_9$ exactly four irreducible curves have zero area, the two nodal curves and
the two exceptional curves of the del Pezzo surface, and they are disjoint $(-1,-1)$-curves
(Proposition~\ref{prop:face}); so the massless spectrum, in all degrees and genera, consists of four
hypermultiplets. No divisor collapses and all gauge couplings stay finite on $F$,
and the resulting threefold with four nodes smooths, in the direction allowed by its single
relation, to a smoothing of $X_9$ deformation equivalent to $X_{9,t}$
(Proposition~\ref{prop:face-smoothing}). The four vanishing charges and their relation were
found from the mirror periods of a four-node degeneration in~\cite[Sects.~5.1--5.2,
Prop.~5.1]{W26quantum}, and the corresponding mirror family was related there to that of
$X_{9,t}$~\cite[Prop.~5.2]{W26quantum}. What is new here is the statement in
five-dimensional M-theory, that on $F$ exactly these curves and no others have zero area, together with a
direct proof that the four-node threefold smooths to $X_{9,t}$. At a root over $F$ the light
states are four free hypermultiplets of a rank-three torus with relation $a=(1,1,-1,1)$, and
the local space of supersymmetric Minkowski vacua is (Theorem~\ref{thm:four-node})
\begin{equation}
 \mathcal V_x\cong(\mathcal K_\phi\times N)\cup(F_\phi\times N\times\CC^2/\ZZ_4),
 \label{eq:intro-main}
\end{equation}
where $\mathcal K_\phi$ is the germ of the Coulomb branch, $F_\phi$ the germ of the face and $N$
the manifold of $123$ neutral hypermultiplets, the two branches being glued along
$F_\phi\times N\times\{0\}$. The transverse type is the rigid one~\eqref{eq:intro-A3}: gravity
does not remove the $A_3$ point or change its type. The leading curvature of the transverse $\CC^2/\ZZ_4$ is the constant curvature $4\nu$ fixed by
supergravity plus a $\ZZ_4$-invariant anti-self-dual Weyl tensor; the latter is the only
compactification-dependent datum at this order. It lies in a space of dimension three, every element of
which comes from some invariant algebraic curvature tensor of hyperk\"ahler type
(Proposition~\ref{prop:x9-leading}); whether each comes from a quaternionic K\"ahler manifold with this
torus action is not addressed. Curvature tensors of hyperk\"ahler type correspond to quartic
polynomials. One of the three directions, invariant under a triholomorphic circle, is realised
by the curvature of the Taub--NUT space at its nut; the other two come from the generators $X$, $Y$
of the ring~\eqref{eq:intro-A3}, which are quartic because $m=4$.

A census of the hypersurfaces from favourable Kreuzer--Skarke polytopes with $h^{1,1}\leq5$ and from the
favourable members of a random sample with $h^{1,1}=6$ (Section~\ref{subsec:census}) finds
5605 faces of toric K\"ahler cones whose light spectrum, computed up to degree six, has a single
relation: 5233 flop walls with $k=2$, and 372 faces with $k=3$, $4$ or $5$ in 210 polytopes. In none of
them does a coefficient $|a_i|\geq2$ occur, so $m=k$ throughout; whether faces with $m\neq k$ exist at all
remains open. For 4092 of the 5605 faces toric data show that the contraction has exactly $k$ nodes, the light
states being disjoint $(-1,-1)$-curves; in the others they do not decide it.

The del Pezzo point, where the whole surface collapses and the $E_2$ theory appears, lies on the
facet of the closure of $F$ on which the del Pezzo surface has zero volume. Call the part of the Higgs branch on which M-theory is compactified
on the smooth threefold $X_{9,t}$ the \emph{regular Higgs branch}. The facet is not a boundary of
it: the regular Higgs branch continues across the facet, and its vacua immediately beyond the facet end
at roots where a surface of a flopped threefold collapses and the $E_1$ theory, together with the two
nodes, appears
(Section~\ref{subsec:other-roots}). Since the $E_2$ theory has no Lagrangian, we define the Higgs
branch at the del Pezzo point as the metric completion of the regular Higgs branch. Because the hypermultiplet metric of
that part does not depend on the vector moduli (Lemma~\ref{lem:product}), the germ computed
over $F$ is carried across the facet, and the completion again has transverse type
$\CC^2/\ZZ_4$ (Theorem~\ref{thm:wall}).

Fourth, in type~IIA string theory on $\widehat X$ in four dimensions, under the corresponding
four-dimensional regularity hypothesis, at large volume, the Wilsonian vacuum space at the locus where the D2-branes on the
curves $C_i$ become massless has the same structure: the hypermultiplet factor of the Higgs branch is
locally $N\times\CC^2/\ZZ_m$, quaternionic K\"ahler off $N\times\{0\}$ and flat up to relative corrections of
order $r^2$ (Theorem~\ref{thm:4d}). The massless states carry no D0-charge, the four-dimensional vacuum
conditions~\cite{Hristov:2009} replace those of five-dimensional supergravity, and at large volume the
special K\"ahler metric is the five-dimensional gauge kinetic matrix divided by $4\operatorname{vol}$
(Proposition~\ref{prop:4d-kinetic}), so the positivity assumed in five dimensions carries over. The $k$
vanishing three-spheres of the smoothing satisfy $[S_i]=a_i[S]$ (Lemma~\ref{lem:spheres}); for unit
coefficients they are homologous up to sign. If in addition the Higgs vacua are type~IIA compactifications on the
smoothing, with hypermultiplet moduli space bi-Lipschitz equivalent to the quotient with its quotient metric, then the type~IIA
hypermultiplet moduli space at finite Planck mass keeps its $A_{k-1}$ point near the $k$ homologous
vanishing three-cycles and is quaternionic K\"ahler around it (Corollary~\ref{cor:gmv}): this is the
statement Greene, Morrison and Vafa expected, for the smoothings of the contractions considered here.

\paragraph{Outline of the proofs.}
The normal form rests on three facts about a fixed point $p$ with vanishing moment map. At a
fixed point the moment map is the $\mathfrak{sp}(1)$-part of the isotropy
(Lemma~\ref{lem:isotropy}), so $\mu(p)=0$ forces the isotropy to preserve a flat hyperk\"ahler
structure on $T_pM$. Along normal geodesics the moment map is the flat one up to terms of order
four, with no cubic term (Lemma~\ref{lem:expansion}), because the fundamental vector field is a
Jacobi field vanishing at the starting point. Blowing up along $N$, the rescaled zero set is
transverse on the link at $r=0$, and the flow of a connection carries the flat cone to the true
zero set; the absence of the cubic term is what makes the deviation of order $r^2$ rather than
$r$. The general physical theorem then reduces to its flat model: for a single relation with
nonzero coefficients the torus acts freely off the origin of the zero set of the flat moment
map, and the flat quotient is $\CC^2/\ZZ_m$ (Lemma~\ref{lem:flat-model}). At the root the
moment map vanishes because the Coulomb branch through it consists of supersymmetric vacua (the
first clause of the hypothesis) and the gauge kinetic matrix is nondegenerate there, so that the
vacuum conditions force every moment map to vanish. For $X_9$ the face is found by a contraction
argument: every $J\in F$ is a class pulled back from an ample class on $X_9$ plus a positive multiple
of a nef class that restricts to the hyperplane class on the del Pezzo surface, so a curve of zero area
is contracted by the resolution and, if it lies in the surface, is one of its exceptional
curves.

\paragraph{Relation to earlier work.}
With gravity decoupled, the $\CC^2/\ZZ_k$ structure of the hypermultiplet moduli space near $k$
homologous vanishing three-cycles was derived from D-instanton sums by Ooguri and
Vafa~\cite{OoguriVafa:1996} and by Seiberg and Shenker~\cite{SeibergShenker:1996}, the
multiplicity was identified with a Gopakumar--Vafa invariant~\cite{SaueressigVandoren:2007}, and
smooth hyperk\"ahler model geometries of this kind were constructed rigorously for $k=1$
in~\cite{FredricksonZimet:2025}, where the extension to multiplicity $k$ is indicated. Greene, Morrison and Vafa obtained $\CC^2/\ZZ_k$ for $k$
hypermultiplets with one relation of unit coefficients~\cite[Sect.~3, eqs.~(3.13)--(3.18)]{Greene:1996Confinement}
and expected that with gravity its qualitative features persist while the metric is
corrected~\cite[Sect.~6]{Greene:1996Confinement}, a statement about the four-dimensional
hypermultiplet space. Theorem~\ref{thm:general} and Proposition~\ref{prop:tangent-cone} make
the five-dimensional counterpart precise, under the regularity hypothesis, for every face with
one relation with nonzero coefficients, and Corollary~\ref{cor:gmv} the four-dimensional statement
itself, for the smoothings considered here and under the identification (W-IIA) of
Section~\ref{subsec:gmv}. Mohaupt and Saueressig constructed five-dimensional
gauged supergravities for M-theory conifold transitions with a model hypermultiplet
metric~\cite{Mohaupt:2004de}; the germ found here does not depend on such a choice. On the mathematical side, quaternionic K\"ahler reduction is due to Galicki and
Lawson~\cite{Galicki:1988} and the correspondence with hyperk\"ahler cones to
Swann~\cite{Swann:1991}. Dancer and Swann described a singular quaternionic K\"ahler quotient as a
union of smooth pieces, some quaternionic K\"ahler and some locally K\"ahler, and as a union of
quaternionic K\"ahler orbifolds~\cite[Thm.~4.6, Props.~4.4 and~5.1, Thm.~6.1]{DancerSwann:1997}, and
showed that for compact connected $M$ and $G$ each component is a complete length space~\cite[Thm.~7.1]{DancerSwann:1997};
through a fixed point with vanishing moment map the piece is the fixed component $N$
(Lemmas~\ref{lem:isotropy} and~\ref{lem:fixed}). These works describe the pieces but not how they fit
together: the frontier condition is not established~\cite[Sect.~2]{DancerSwann:1997}, and no local
model near a singular piece is given. In symplectic geometry that role is
played by the local normal form of Marle and of Guillemin and
Sternberg~\cite{Marle:1985,GuilleminSternberg:1984}, on which the stratification of singular symplectic
quotients by Sjamaar and Lerman~\cite{SjamaarLerman:1991} rests. The hyperk\"ahler analogue of that
local model is not a normal form, since a hyperk\"ahler equivalence is an
isometry~\cite[p.~149]{Swann:1997sing}; see
also~\cite[Sect.~1]{Mayrand:2022}. For hyperk\"ahler quotients of positive-definite hyperk\"ahler
manifolds on which the complexified group acts, that is, on which the vector fields $I_1X_\xi$ are
complete, Mayrand gives a complex-analytic local model at every point of the
quotient~\cite[Thm.~1.4]{Mayrand:2022}, and the product form of the germ along a fixed component was
conjectured for semi-free circle actions by Fang~\cite[Rem.~3.2]{Fang:2008}. For negative scalar
curvature the Swann bundle is pseudo-hyperk\"ahler, so the model of~\cite{Mayrand:2022} does not
apply, and Theorem~\ref{thm:normal-form} is proved directly on the quaternionic K\"ahler manifold;
for positive scalar curvature that model describes the complex-analytic germ of the hyperk\"ahler
cone, which is not the metric statement along $N$ proved here. Theorem~\ref{thm:normal-form} is
accordingly a model up to homeomorphism of germs, with metric estimates, and
Proposition~\ref{prop:tangent-cone} describes the tangent cone of the orbit length metric at such
points without assuming compactness; for compact connected $M$ and $G$ this metric is that of Dancer and Swann.

\paragraph{Scope.}
The results of Section~\ref{sec:quotients}, among them Theorem~\ref{thm:normal-form} and
Proposition~\ref{prop:universal}, are theorems about quaternionic K\"ahler quotients and are proved
unconditionally, as are Propositions~\ref{prop:face}, \ref{prop:face-smoothing}
and~\ref{prop:4d-kinetic} and Lemmas~\ref{lem:flat-model} and~\ref{lem:spheres}.
Lemmas~\ref{lem:product} and~\ref{lem:completion} and Theorems~\ref{thm:general},
\ref{thm:four-node} and~\ref{thm:wall} are proved under Hypothesis~\ref{hyp:regular};
Theorem~\ref{thm:4d} and Corollary~\ref{cor:gmv} under Hypothesis~\ref{hyp:regular-4d} and, for the
corollary, the identification (W-IIA) of Section~\ref{subsec:gmv}. Proposition~\ref{prop:x9-leading} is proved for the manifold $M$ of Hypothesis~\ref{hyp:regular};
Appendix~\ref{app:data} describes an independent computer check of it and of
Proposition~\ref{prop:universal}. The
census of Section~\ref{subsec:census} rests on computer calculations of genus-zero invariants. Theorem~\ref{thm:wall} describes the metric completion of the
part of the Higgs branch over the smooth threefold $X_{9,t}$, over a region of the K\"ahler cone of $X_{9,t}$ that contains the del
Pezzo facet, not the vacuum space of the interacting theory at the del Pezzo point. The
statements about the other roots of $X_9$ in Section~\ref{sec:discussion} concern their light
spectra and rigid models only. All statements about the vacuum space are topological or metric;
no complex-analytic structure is claimed.

\paragraph{Organisation.}
Section~\ref{sec:face} recalls $X_9$ and proves the statements about the face.
Section~\ref{sec:vacua} states the setting for a general face with one relation, the vacuum
conditions and the regularity hypothesis. Section~\ref{sec:quotients} proves the normal form and
the tangent-cone statement and identifies the leading curvature. Section~\ref{sec:higgs} proves the theorems on the Higgs branch and computes the leading correction
for $X_9$. Section~\ref{sec:4d} treats type~IIA string theory in four dimensions.
Section~\ref{sec:discussion} compares with the rigid picture, describes the other roots of $X_9$, compares
with the instanton picture, reports the census of one-relation faces, discusses other effects at finite
Planck mass, and lists open problems. Appendix~\ref{app:data} prints the intersection data, the kinetic
determinant, the Gopakumar--Vafa data and the charge conventions, and describes an independent check of
the leading correction.
